% P1 one-step policies -> P2 the fine-tuning gap -> P3 toy/flaw -> P4 key idea
% -> P5 dial + results -> contributions.

A behavior-cloned robot policy is only as good as its demonstrations. To
push past them, the usual move is to learn a critic and fine-tune against
it with reinforcement learning. This paper is about one design choice
inside that loop, which we think is made wrongly almost everywhere: how
much authority the critic is given.

A critic does two jobs in a standard update. It says which way the policy
should move, and --- through the magnitude of its gradient --- how far. Only
the first is a job it was trained for. A critic learns to rank actions in
the region its data covers; outside that region nothing constrains it, and
it is quite capable of being steep and wrong at the same time
\citep{fujimoto2018td3,fujimoto2019bcq,savo2025}. The consequences of the
two errors are not comparable. A slightly wrong direction gives a slightly
wrong step. A wrongly large distance takes the policy off the data manifold
in one update, and the further off it goes the better it looks to the
critic that sent it there.

Our setting sharpens the problem, because it removes the usual escape
route. We fine-tune \emph{one-step} generative policies --- a drifting
model \citep{deng2026drifting} maps noise to a multimodal action in a
single network evaluation --- and such a policy has no tractable
likelihood, so the ratio-based toolbox is unavailable. It can be recovered
only by injecting noise into the deployed policy \citep{reinflow2025},
which changes what ships, or by assuming a denoising objective the policy
does not have \citep{mcallister2025flowmatchingpolicygradients}
(\cref{sec:failure-modes}). With likelihoods off the table, a critic is
what remains --- which makes how we consume it the whole question.

Both standard answers delegate the distance. Backpropagating $-Q$ through
the actor, as the state-of-the-art residual fine-tuner DICE-RL
\citep{dicerl2026} does, makes every parameter step proportional to
$\|\nabla_a Q\|$. The action-gradient family
\citep{dipo2024,psenka2024qsm} never differentiates through the policy, but
moves each sampled action to wherever a gradient step lands and regresses
onto that --- the same delegation, one level up. A controlled toy
(\cref{sec:toy}) shows how bad this can get: with high-dimensional actions
and scarce critic data, the backpropagated update drives true reward to
zero while its own $Q$-estimate rises by four orders of magnitude. A
behavior-cloning penalty does not save it, because a penalty in the loss
must out-shout an unbounded gradient, and loses. Yet the gradient is not
the problem. In the opposite regime, where the critic generalizes, it is
the best signal available, and methods that discard it waste it. What is
needed is a rule that keeps the direction in both regimes while capping how
far any single step may carry the policy.

The lesson we draw is narrow and we state it as such: \emph{a critic may be
trusted with the direction of a policy update, and with its length only up to
a bound}. We are not claiming the critic should be stripped of the step size
entirely --- in our own method it still sets it for the majority of sampled
actions (\cref{sec:exp-ablation}), and that is fine. The claim is that the
\emph{tail} must be truncated, because that is where the critic is least
constrained and where a single step can end the run.
Discarding the direction, as critic-as-judge methods do
\citep{peng2019awr,wang2020crr}, throws away the best signal available when
the critic is trustworthy; letting the critic set the length hands it a
quantity it was never trained to supply. Note that this indicts more than
backpropagation. The \emph{action-gradient} family --- which improves
actions along $\nabla_a Q$ and regresses the policy onto the improved
actions \citep{dipo2024,psenka2024qsm} --- keeps the direction but still
leaves the distance to the critic, and is reported to degrade exactly as
one would predict, drifting off the behavior distribution when the critic's
gradient is biased \citep{dprlsurvey2026}. What has been missing across
both is not a better signal but a \emph{bound}, and one placed where the
harm occurs: in action space, not parameter space. Gradient clipping
constrains the parameter step, which says nothing about how far the
resulting action moves.

\method{} supplies that bound. A frozen base proposes action particles
through a residual head, and reward enters as a velocity field $V_Q$ --- by
default the normalized critic gradient. Two guards then act before anything
is regressed. First, each particle's total displacement is \emph{clipped
pointwise in action space}, so the policy provably moves at most $\delta$
per update at every particle regardless of critic pathology --- a guarantee
parameter-space gradient clipping cannot give, since a bounded parameter
step says nothing about the distance the action travels. Second, a
restoring anchor field with a dead zone pulls the residual back toward the
base whenever it strays beyond a radius $\rho$, leaving it untouched
inside. Only then is the residual regressed onto the displaced targets, and
because the update rewrites nothing in the replay buffer, stored
transitions remain consistent with the dynamics that produced them. No
likelihood is ever needed, and in the small-step limit the update recovers
the backpropagated direction exactly (\cref{prop:equiv}) --- so the safety
costs nothing asymptotically.

The construction is also economical on the policy class that motivated it.
A drifting model \citep{deng2026drifting} is \emph{pretrained} by
field-target regression --- transport samples along a velocity field,
regress the network onto the transported targets --- so on that base the
fine-tuner reuses machinery the policy already has. That economy is a
convenience, not a requirement: the update assumes only a frozen
noise-to-action map, and we verify it transfers unchanged to a multi-step
flow-matching base.

Because reward enters only as a field, the critic can also be consumed at
lower resolution when its gradients cannot be trusted: transport toward
critic-ranked top-$k$ particles, or toward the cloud's own
$\exp(Q/\tau)$-tilted density. Our experiments treat this as an explicit
dial and map its regimes. On robomimic, where the critic generalizes, the
gradient field is the resolution to use: from a drifting base at $0.877$, \method{}
reaches $0.99$ on can, and lifts square from $0.382$ to
$0.91$--$0.93$ across seeds --- ahead of the backpropagated control at the
matched budget ($0.87$--$0.91$) and at parity with the flow-matching +
DICE-RL state-of-the-art pipeline on both tasks, while the value-only
sources trail far behind ($\sim\!0.55$ on square). On the eight
hardest LIBERO-90 tasks, fine-tuned per task from one shared base with
one shared recipe, \method{} gains $+3.9$ to $+23.9$ points (mean
$+9.4$) where the
backpropagated actor consuming the same critics gains $+0.9$ and
\emph{loses} ground on the two highest-headroom tasks. The ordering is
unchanged when the base is swapped for the stronger flow-matching policy,
and \method{} matches a critic-free policy-gradient baseline at equal
budget without needing the noise injection that baseline requires. With a
single critic shared across all eight tasks, no \emph{critic-based} update
rule --- backpropagated or field-based, ours or the state of the art's ---
improves on its own base, locating the binding constraint in per-task
critic quality exactly as the toy predicts.

Our contributions:
\begin{itemize}
  \item \textbf{A dead-zone restoring force for action-gradient policy
  improvement} (\cref{sec:method}). Improving actions along $\nabla_a Q$ and
  regressing the policy onto them is an established family \citep{dipo2024,
  psenka2024qsm, dprlsurvey2026}, and we adopt it unchanged. What that family
  lacks is anything holding the policy near the data, and the obvious
  candidate --- a behaviour-cloning penalty applied at every step --- is the
  wrong shape, because it opposes the critic in the region where the critic is
  reliable. We instead apply a restoring pull that is \emph{silent} within a
  radius $\rho$ of the frozen base and engages only beyond it, matching the
  regularizer's spatial profile to the critic's: free rein where the critic
  has data, a boundary where it does not. The update needs no likelihood.
  \item \textbf{A controlled diagnosis of critic exploitation}
  (\cref{sec:toy}): a two-regime toy showing the failure is real,
  regime-dependent, and \emph{not} fixed by BC penalties --- motivating
  trust regions in action space rather than loss space.
  \item \textbf{Evidence that the \emph{form}, not the signal or the step
  size, is what wins} (\cref{sec:experiments}): holding base, critic, data
  and budget fixed and varying only the actor update, \method{} beats the
  backpropagated actor on both benchmarks and on both bases, lifts a weaker
  one-step base to match the stronger flow-matching + DICE-RL pipeline, and
  reaches parity with a critic-free policy-gradient baseline without
  modifying the deployed policy. Two controls close the obvious
  alternative explanations. Sweeping each method's own step size an order
  of magnitude leaves the two in non-overlapping bands --- the
  backpropagated actor never rises above its base at any learning rate,
  ours never falls near it at any reward step --- so the gap is not
  tuning. And deleting only the critic's gradient direction, keeping
  everything else, collapses performance \emph{below} the untuned base on
  two separate tasks --- so the gain is not simply caution.
\end{itemize}
